\documentclass{article}
\usepackage[T1]{fontenc}
\usepackage[style=authoryear, backend=biber]{biblatex} %load package
\addbibresource{references.bib} %points to bib file
\usepackage{array} %tables
\usepackage{booktabs} %line rules

\title{testinglatex}
\author{me}
\date{\today}

\begin{document}

\maketitle 
trying to lock in and learn latex xx
\bigskip

emph is context-aware. \emph{it makes text \emph{stand out} from whatever text surrounds it}.

textit is a font-switch. it \textit{blindly forces} whatever is in it to be italic

textbf is also a font-switch. it \textbf{converts text to boldface}.

\section{title of the section}
some text
\subsection{title of the subsection}
some text
\bigskip

ordered list
\begin{enumerate}
\item an entry
\item ooh another one
\item 3rd one
\end{enumerate}

unordered list
\begin{itemize}
\item an entry
\item 2nd one
\item wow this is fun
\end{itemize}

this is a sample document with some dummy text\footnote{and a footnote}.

\section{tables}

\subsection{simple tables}

\begin{tabular}{lll}
flower & color & price \\
rose & red & 2 \\
peony & pink & 3 \\
\end{tabular}

\subsection{text-wrap tables}

\begin{tabular}{cp{9cm}}
    \toprule
    something & something \\
    \midrule
    hi & this format is basically for long text. if we use the previous format, we end up with the text running off the paper\\
    \cmidrule{1-2}
    bye & also makes it look super clean and i dont know what to type anymore- my stomach hurts and i miss my dog so yes ok bye\\
    \cmidrule{1-2}
    wait & oh also i added rules bc i think it looks better. {1-2} gives an complete line, {1-1} for the first column, {2-2} for the second column\\
    \bottomrule
\end{tabular}

\section{math}

two forms of math mode: inline and display. inline is used for eqns part of a sentence, display is pretty math
like, sentence with inline mathematics: $y = mx + c$.
another sentence with inline mathematics: $5^{2}=3^{2}+4^{2}$.
\bigskip

PRETTY MATHHH
\[
  y = mx + c
\]

superscripts $x^{y}$ and subscripts $x_{y}$.
\medskip

trig: $y = 2 \sin \theta^{2}$.
\medskip

\begin{equation}
    \int_{-\infty}^{+\infty} e^{-x^2} \, dx
\end{equation}

note so i don't lose my mind later: the asmath package has more stuff

\section{citing sources}

according to cosmology, time has a beginning \cite{hawking1988}. 
albert einstein also made massive contributions to physics \cite{einstein1905}.
\bigskip
\bigskip

okay i think this is enough for now i hope i dont need anything more. YAY latex!! time to sleep byeeeeeeeeee

\printbibliography

\end{document}